\section{Literature Review}\label{sec:literature}

This section reviews five lines of prior work and then states the research gap (Section~\ref{sec:lit-gap}).

\subsection{Misconfigurations in human-written Kubernetes configurations}\label{sec:lit-human}

Rahman et al.\ identified seven security smells in Puppet scripts, such as hard-coded secrets \citep{rahman2019sevensins}. For Kubernetes, Shamim et al.\ organised practitioner knowledge into eleven recommended security practices \citep{shamim2020xicommandments}. The most complete study of Kubernetes security misconfigurations analysed 2,039 manifests from 92 open-source repositories. Its 11 categories, such as absent resource limits and privilege escalation, affected 13.9\% of the analysed objects \citep{rahman2023k8smisconfig}. Beyond security, \citet{zhang2026k8sconfigdefects} derived 15 defect categories from 719 configuration defects and showed that eight static analysis tools could detect only 8 of them. For Helm charts, the most frequent misconfigurations concerned resource quotas, network policies and security profiles \citep{minna2025helm}, and only 21.43\% of the misconfigurations in incubator charts were fixed by the stable stage \citep{ghorab2026inthewild}.

We reuse their taxonomies, but these studies observe only human-written configurations, and nearly all analyse them statically without deploying them.

\subsection{Standards and tools for assessing Kubernetes configurations}\label{sec:lit-tools}

Guidelines such as the NSA/CISA Kubernetes Hardening Guide \citep{nsacisa2022k8shardening}, the CIS Kubernetes Benchmark \citep{cis_k8s_benchmark}, the OWASP Kubernetes Top Ten \citep{owasp_k8s_top10} and the Kubernetes Pod Security Standards \citep{k8s_pod_security_standards} define a secure configuration, and static scanners put them into practice. Kubescape \citep{tool_kubescape}, Trivy \citep{tool_trivy} and Checkov \citep{tool_checkov} focus mainly on security. kube-score \citep{tool_kubescore}, Polaris \citep{tool_polaris} and KubeLinter \citep{tool_kubelinter} also check reliability practices, such as probes and resource requests. kubeconform checks manifests against the official Kubernetes API schemas \citep{tool_kubeconform}.

These tools often disagree. On a benchmark of 79 recommendations, only two of ten scanners reached an F1-score of at least 0.69 while covering at least half of the recommendations, and severity ratings were inconsistent \citep{krieger2026k8sscanners}. \citet{kapetanidou2025k8sscanning} recommended adding in-cluster workload scanning to static analysis, and \citet{lanciano2023declarative} used Polaris findings as quality labels. So one scanner alone risks tool-specific bias, and static verdicts need runtime evidence.

\subsection{Security of LLM-generated code}\label{sec:lit-llmcode}

LLM-generated code often has security flaws. About 40\% of the programs GitHub Copilot wrote for CWE Top-25 scenarios were vulnerable \citep{pearce2022asleep}. Similar weaknesses appear in programs written by ChatGPT \citep{khoury2023chatgptsecure}, in AI-generated code committed to GitHub \citep{fu2025copilotweaknesses} and in about 30\% of the CyberSecEval test cases \citep{bhatt2023cyberseceval}. ChatGPT could explain the vulnerabilities when asked but did not apply this knowledge unprompted \citep{khoury2023chatgptsecure}. BaxBench showed that about half of the functionally correct backends were exploitable \citep{vero2025baxbench}.

So functional correctness does not imply security, and LLMs tend to omit security measures unless asked. But these studies cover imperative code. In declarative configurations, security and reliability depend largely on \emph{what is omitted}, such as a security context, probe or resource limit, and this failure has received far less attention.

\subsection{LLM-based generation of infrastructure and cloud configurations}\label{sec:lit-iac}

Ansible Wisdom fine-tuned a code model on Ansible YAML \citep{pujar2023ansiblewisdom}. CloudEval-YAML has 1,011 Kubernetes, Envoy and Istio problems, each with a unit test run on a real cluster; the best of 12 models (GPT-4) scored only 0.515 \citep{xu2024cloudevalyaml}. The best model reached a pass@1 of only 19.36\% on the Terraform benchmark IaC-Eval \citep{kon2024iaceval}, and no open-source model exceeded 12\% pass@10 on the Ansible benchmark ITAB \citep{hassan2026itab}. TerraFormer fine-tuned models with feedback from formal verifiers and policy checks \citep{jana2026terraformer}, and \citet{kozachok2026k8sdistillation} used staged validators, including kubeconform, to adapt small language models to manifest generation.

Recent studies also measure deployability and security. Only 20.8--30.2\% of CloudFormation templates from six LLMs deployed at the first attempt, and security compliance stayed at 8.4\% even after iterative refinement \citep{zhang2026iacgen}. For Terraform, syntactic validity and security compliance were ``largely orthogonal'' \citep{santosvargas2026secureterraform}. FDE-Bench evaluated LLM agents on 136 deployment-configuration tasks, including Kubernetes, and readiness was the main failure stage \citep{ding2026fdebench}.

So validity, deployability and security diverge, and each must be measured. But these works use curated or synthetic tasks with reference solutions. None compares LLM outputs with the configurations developers maintain for real applications, or tests behaviour under failures.

\begin{table*}[t]
\caption{Comparison of the most closely related studies with this work. Gen.\ = generation from requirements; Det./Rep.\ = detection or repair of existing configurations; \checkmark\ = covered; $\circ$ = partially covered; -- = not covered.}\label{tab:related}
\centering
\resizebox{\textwidth}{!}{%
\begin{tabular}{@{}llllcccccc@{}}
\toprule
Study & Artefact & Task & LLMs & \shortstack{Matched\\human} & \shortstack{Multi-\\tool} & Security & Reliability & Deploy & Recovery \\
\midrule
\citet{rahman2023k8smisconfig} & K8s & -- (human only) & -- & -- & -- & \checkmark & $\circ$ & -- & -- \\
\citet{zhang2026k8sconfigdefects} & K8s & -- (human only) & -- & -- & \checkmark & $\circ$ & \checkmark & -- & -- \\
\citet{xu2024cloudevalyaml} & K8s/Istio/Envoy & Gen. & 12 & -- & -- & -- & -- & \checkmark & -- \\
\citet{zhang2026iacgen} & CloudFormation & Gen. & 6 & -- & -- & $\circ$ & -- & \checkmark & -- \\
\citet{santosvargas2026secureterraform} & Terraform & Gen. & 7 & -- & \checkmark & \checkmark & -- & -- & -- \\
\citet{ding2026fdebench} & Docker/K8s & Gen. (agentic) & 7 & -- & -- & -- & $\circ$ & \checkmark & -- \\
\citet{zhang2024genaismells} & K8s & Gen. (in the wild) & 1 & -- & -- & \checkmark & $\circ$ & -- & -- \\
\citet{jcomsec2026aivshumank8s} & K8s & Gen. (reused) & 1 & -- & -- & \checkmark & $\circ$ & -- & -- \\
\citet{malul2024genkubesec} & K8s & Det./Rep. & 5+ & -- & \checkmark & \checkmark & -- & -- & -- \\
\citet{ye2025llmsecconfig} & K8s/Helm & Rep. & 2 & -- & -- & \checkmark & -- & -- & -- \\
\citet{sgancohen2025kubeguard} & K8s & Hardening & 2+ & -- & -- & \checkmark & -- & $\circ$ & -- \\
\midrule
\textbf{This study} & K8s & Gen. & 5 & \checkmark & \checkmark & \checkmark & \checkmark & \checkmark & \checkmark \\
\bottomrule
\end{tabular}}
\end{table*}

\subsection{Studies of LLM-generated Kubernetes manifests}\label{sec:lit-llmk8s}

Two studies examine LLM-generated Kubernetes manifests directly. \citet{zhang2024genaismells} analysed 98 ChatGPT-generated manifests from the DevGPT dataset with KubeLinter. Unset CPU and memory requirements were the most frequent smells, and each affected 35.8\% of the manifests. The authors noted that their findings were limited to one model, one tool and manifests of unknown origin. \citet{jcomsec2026aivshumank8s} compared these measurements with 1,985 human-written manifests from the corpus of Rahman et al. Missing CPU and memory specifications were the most robust weakness of the AI-generated manifests, while missing \texttt{runAsNonRoot} and \texttt{readOnlyRootFilesystem} settings were common in both corpora.

These studies show that LLM-generated manifests inherit known weaknesses and sometimes amplify them. The authors note the limits: the corpora are unmatched and differ greatly in size, the AI corpus reflects a single 2023 model, the analysis uses one tool, and no manifest is run in a cluster \citep{jcomsec2026aivshumank8s}.

\subsection{LLMs for detecting, repairing and hardening Kubernetes configurations}\label{sec:lit-llmsec}

Other work uses LLMs on existing configurations. Fine-tuned models matched or exceeded rule-based tools in detecting misconfigurations \citep{malul2024genkubesec,ghorab2025k8sllm}, while off-the-shelf models such as ChatGPT, Gemini and Claude did poorly (best F1 $\approx$ 0.51) \citep{malul2024genkubesec}. For repair, LLMSecConfig combined Checkov findings with retrieval-augmented generation and found large differences in repair success between models (94.3\% for Mistral Large 2 versus 40.2\% for GPT-4o-mini) \citep{ye2025llmsecconfig}. Schema-guided correction raised LLM repair accuracy to 98.5\% \citep{ghorab2026inthewild}. Yet LLM fixes for Helm charts can satisfy scanners while deleting unrelated configuration and breaking the application \citep{minna2025helm}. KubeGuard derived least-privilege roles, network policies and deployments from runtime logs \citep{sgancohen2025kubeguard}. KuTIE showed that LLM security patches that ignore the live service topology are correct in only 11.1\% of topology-dependent cases \citep{shaikh2026kutie}.

So models disagree, and a scanner-compliant configuration can still break a running service. But all these works start from an existing manifest. None asks how well an LLM writes the complete configuration of an application from a requirements description, as developers already do \citep{zhang2024genaismells}.

\subsection{Research gap}\label{sec:lit-gap}

Table~\ref{tab:related} compares the closest studies with ours. The review shows five gaps.
\begin{enumerate}
    \item \textbf{Real-project comparison.} Benchmarks use curated tasks, and the only human-versus-AI Kubernetes comparison uses unmatched corpora. No study compares LLM and human configurations of the \emph{same} applications, generated from the same requirements.
    \item \textbf{Current LLMs.} Evidence rests on a single 2023 ChatGPT corpus. Current proprietary and open-weight models, and their consistency across repeated generations, have not been systematically compared.
    \item \textbf{Reliability.} Work on LLM-generated configurations focuses on security. Reliability controls (probes, resource requests, replicas, disruption budgets, anti-affinity) have been studied only for human-written configurations.
    \item \textbf{Runtime resilience.} Prior work stops at static analysis or a ready state. No study injects failures to test whether generated configurations recover and keep the service available.
    \item \textbf{Linking static and runtime evidence.} Scanners disagree, and compliant configurations can still fail at runtime. No study evaluates the same configurations on static correctness, security, reliability, deployment and recovery to test whether static findings predict runtime behaviour.
\end{enumerate}
This study addresses these gaps by comparing human-written configurations with those of five current LLMs for the same projects on a six-dimension benchmark that covers static analysis, requirement satisfaction, real deployment and fault injection.
